% Topologie L2/L3 — vue logique
\begin{tikzpicture}[node distance=0.85cm]
  \node[term, minimum width=3.6cm] (inet) {Internet};
  \node[blk, below=of inet, minimum width=3.6cm] (isp) {Opérateur fibre};
  \node[key, below=of isp, minimum width=5.2cm, minimum height=1.3cm] (srx)
    {\textbf{Juniper SRX300}\\[1pt]\scriptsize NAT · VPN · filtrage};
  \node[blk, below left=1.1cm and 0.6cm of srx, minimum width=3.8cm] (sw1) {Switch Cisco 1};
  \node[blk, below right=1.1cm and 0.6cm of srx, minimum width=3.8cm] (sw2) {Switch Cisco 2};
  \node[blk, below=1.5cm of sw1, minimum width=4.4cm] (r630) {R630 $\times$ 4\\[1pt]\scriptsize Proxmox (transitoire)};
  \node[blk, below=1.5cm of sw2, minimum width=4.4cm] (hist) {Serveurs historiques\\[1pt]\scriptsize cible OpenStack};

  \draw[flow] (inet) -- (isp);
  \draw[flow] (isp) -- node[lab, right=2pt] {fibre} (srx);
  \draw[flow] (srx.south) -- ++(0,-0.45) -| (sw1.north);
  \draw[flow] (srx.south) -- ++(0,-0.45) -| (sw2.north);
  \draw[dflow] (sw1) -- node[lab, above=1pt] {trunk LACP} (sw2);
  \draw[flow] (sw1) -- node[lab, left=3pt] {VLAN 10 · 20 · 40} (r630);
  \draw[flow] (sw2) -- node[lab, right=3pt] {VLAN 10 · 20 · 30} (hist);
\end{tikzpicture}
